% ================================================================
% 02_more_theorem.tex — 更多定理环境帧
% 覆盖: 引理/命题/推论/注记 + 其余全部编号环境 (全量测试)
% 注意: Beamer 帧内容不可跨页, 环境较多时拆成多帧。
% ================================================================

\begin{frame}{更多定理环境}
  \begin{lemma}{}
    若 $f$ 在闭区间上连续, 则 $f$ 一致连续。
  \end{lemma}

  \begin{proposition}{}
    连续函数把紧集映射为紧集。
  \end{proposition}

  \begin{corollary}{}
    闭区间上的连续函数必有最大值与最小值。
  \end{corollary}
\end{frame}

\begin{frame}{公理与假设}
  \begin{axiom}{选择公理}{}
    任意非空集合族 $\{A_i\}_{i \in I}$ 存在选择函数。
  \end{axiom}

  \begin{hypothesis}{}{}
    假设 $f$ 在 $[a, b]$ 上 Riemann 可积。
  \end{hypothesis}
\end{frame}

\begin{frame}{元定理与准则}
  \begin{metatheorem}{}{}
    元定理: 所有定理均可被形式化验证。
  \end{metatheorem}

  \begin{criteria}{Cauchy 准则}{}
    数列收敛当且仅当它是 Cauchy 列。
  \end{criteria}
\end{frame}

\begin{frame}{约定与警示}
  \begin{convention}{}{}
    本书中 $\N$ 含 $0$。
  \end{convention}

  \begin{algorithm}{Euclid 算法}{}
    输入 $a, b$; 当 $b \neq 0$ 时令 $(a, b) \leftarrow (b, a \bmod b)$。
  \end{algorithm}

  \begin{alarm}{常见错误}{}
    不要把 $\lim (a_n b_n) = \lim a_n \cdot \lim b_n$ 用于发散数列。
  \end{alarm}
\end{frame}

\begin{frame}{问题与回答}
  \begin{problem}{}{pr:continuity}
    判断 $f(x) = |x|$ 在 $x = 0$ 处是否可导。
  \end{problem}

  \pause

  \begin{answer}
    不可导: 左右导数分别为 $-1$ 与 $1$, 不相等。
  \end{answer}
\end{frame}

\begin{frame}{例、注记与分析}
  \begin{example}{}{}
    $f(x) = x^2$ 在 $x = 0$ 处可导, 且 $f'(0) = 0$。
  \end{example}

  \begin{remark}{}
    上述结论在度量空间框架下仍然成立。
  \end{remark}

  \begin{sketch}{}{}
    由导数的定义直接计算差商极限。
  \end{sketch}
\end{frame}
